% Input fragment. Requires amsmath, amssymb, amsthm and the parent theorem environments.
% The mathematical adaptation was constructed by the alternate-analysis agent;
% it is not an independent review of its own contribution.
\section{A single-coin entropy bound with all-negative reference cost}

Fix $a=1/4$ and $\mu(dd)=2d a^{-2}\mathbf1_{(0,a]}(d)\,dd$.
Throughout this section,
\[
 \mathcal X_\rho=\left\{q:0\le q\le\tfrac12,
                    \ \int q\,d\mu\le\rho\right\},
 \qquad 0<\rho\le\tfrac14,
 \qquad L=1+\log(1/\rho),\quad H=1+\log L.
\]
All functions defining populations are measurable. A single genuine coin
has hidden $D\sim\mu$ and $S\mid D=d\sim\operatorname{Ber}(q(d))$;
conditionally on $(D,S)$ its observations are iid Bernoulli with parameter
$1/2+(2S-1)D$. A policy receives only its independent seed and observed
bits. It may know $\mu,a,\rho$, but not $q,D,S$.

\begin{theorem}[The critical product with null-population cost]
\label{thm:base-cost}
There is a universal $c_0>0$ such that the following holds for every
$0<\rho\le1/4$. Let an observable single-coin procedure have real output
$Y$ and stopping time $T\in\mathbb N_0\cup\{\infty\}$, and suppose
\[
 \sup_{q\in\mathcal X_\rho}|\mathbb E_qY-\theta(q)|\le\rho/64,
 \qquad
 V:=\sup_{q\in\mathcal X_\rho}\operatorname{Var}_qY<\infty,
 \qquad
 \sup_{q\in\mathcal X_\rho}\mathbb E_qT<\infty,
 \qquad \theta(q)=\int q\,d\mu.
\]
Writing $W_0=\mathbb E_{q=0}T$, one has
\[
                 V W_0\ge c_0\rho L H.                 \tag{BC1}
\]
The same order lower bound holds with
$M=\sup_{q\in\mathcal X_\rho}\mathbb E_qY^2$ in place of $V$.
The cost is the expectation under the all-negative population, rather
than the supremum of expected work over the parameter class. The policy
may stop before observing the coin and may have unbounded stopping.
\end{theorem}

The proof adapts the separately documented critical-two single-coin
argument. We reproduce its finite-gap dependency to make this statement
self-contained. Its entropy mechanism follows Chen--Li--Qiao,
arXiv:1608.06031, Lemma~4.1; the Bernoulli, one-sided-cost statement used
here is proved explicitly below. No novelty is claimed for that mechanism.

\begin{lemma}[Finite-gap sign entropy]\label{lem:finite-sign-entropy}
There are universal $c_s>0$ and $K_s<\infty$ with the following property.
Let $d_1>\cdots>d_K>0$ satisfy $d_j/d_{j+1}\ge2$ and
$d_j\le1/128$. Suppose one observable sign-testing procedure, the same
at all these gaps, has error at most $1/10$ under every Bernoulli law
$1/2\pm d_j$, and finite expected stopping time $\tau$ under both signs.
For $K\ge K_s$,
\[
              \sum_{j=1}^K d_j^2\mathbb E_{-d_j}\tau
                   \ge c_s K\log K.                    \tag{BC2}
\]
Correctness and termination under the fair coin are not required.
One may take $c_s=1/512$ and $K_s=\lceil e^{32}\rceil$.
\end{lemma}
\begin{proof}
Put $b_0=1/1024$, $n_j=\lfloor b_0/d_j^2\rfloor$,
$u_j=\lceil4\mathbb E_{-d_j}\tau\rceil$,
$c_j=d_j^2\mathbb E_{-d_j}\tau$, and
$\bar c=K^{-1}\sum_jc_j$. Since $d_j\le1/128$,
$n_j\ge b_0/(2d_j^2)\ge1$. The event of returning minus by time $n_j$
is measurable in the independent seed and the first $n_j$ flips, including
unused coordinates if the test has stopped. The KL divergence between
these full prefixes at the minus and plus laws is at most
\[
 n_j\frac{(2d_j)^2}{(1/2+d_j)(1/2-d_j)}
 \le\frac{64}{3}n_jd_j^2\le\frac1{48}.
\]
Bernoulli KL is at most its chi-square divergence. Pinsker's inequality
bounds the total variation by less than $1/8$. Since an early minus is
an error under the plus law, it has probability at most $1/10+1/8$ under
the minus law. Markov's inequality gives
$\mathbb P_{-d_j}\{\tau>u_j\}\le1/4$. Consequently
\[
 A_j=\{\text{output minus and }n_j<\tau\le u_j\}
 \quad\text{satisfies}\quad
 \mathbb P_{-d_j}(A_j)
 \ge1-\tfrac1{10}-(\tfrac1{10}+\tfrac18)-\tfrac14>\tfrac14.
\]

Let $P_0$ denote the fair-coin law with the same seed and policy. Comparing
the deterministic prefixes of length $u_j$, their KL divergence from
minus to fair is at most
$4d_j^2u_j\le16c_j+1$. For an event of probability $p\ge1/4$ under
the minus law and probability $z$ under the fair law, binary data
processing and $\operatorname{kl}(p,z)\ge p\log(1/z)-\log2$ therefore give
\[
                     P_0(A_j)\ge e^{-64c_j-8}.          \tag{BC3}
\]
Only deterministic finite prefixes have been compared; this does not
invoke optional stopping under a possibly nonterminating fair law.

At least $K/2$ indices form a set $G$ with $c_j\le2\bar c$. For them,
\[
 (n_j,u_j]\subset
 [b_0/(2d_j^2),(8\bar c+1)/d_j^2].
\]
The scales $d_j^{-2}$ have successive ratios at least four. At any finite
time the number of these intervals containing that time is at most
$1+\log_4((16\bar c+2)/b_0)$. All events $A_j$ refer to the same terminal
time and output on a realization under the fair law; on a nonterminating
path none occurs. Summing their indicators and using (BC3) yields
\[
 (K/2)e^{-128\bar c-8}
        \le1+\log_4(1024(16\bar c+2)).                  \tag{BC4}
\]
If $x=\log K\ge32$ and $\bar c<x/512$, the left side exceeds
$\tfrac12e^{3x/4-8}\ge\tfrac12e^{x/2}\ge x^2/16$.
The right side is less than
$1+\log_4(32x+2048)<6+x/4$, contradicting $x\ge32$.
Thus $\bar c\ge(\log K)/512$, as claimed.
\end{proof}

\begin{proof}[Proof of Theorem~\ref{thm:base-cost}]
The bias bound implies $M\le V+(65\rho/64)^2<\infty$.
First define the nonnegative extended conditional moments
\[
 M_s(d)=\mathbb E_{d,s}Y^2,\qquad w_s(d)=\mathbb E_{d,s}T.
\]
The admissible populations $q=0$ and $q\equiv\rho$ imply
$\int M_-d\mu\le M$, $\int M_+d\mu\le M/\rho$, and finite
integrals of both work functions. All conditional quantities used below
are thus finite outside a common null set. Define signed conditional
means there, and give them arbitrary values on that null set. For
$f(d)=\mathbb E_{d,+}Y-\mathbb E_{d,-}Y$, the two bias bounds imply
\[
 \int f\,d\mu=(\mathbb E_\rho Y-\mathbb E_0Y)/\rho\ge31/32.
\]
Conditioning on the latent magnitude and then on its sign at the
constant-prior population gives
\[
 V\ge\operatorname{Var}_\rho(Y)
 \ge\rho(1-\rho)\int f^2\,d\mu
 \ge\rho(1-\rho)(31/32)^2\ge2\rho/3.                   \tag{BC5}
\]
This uses only the bias guarantee, not a pointwise unbiased response.
It also shows
\[
                        M\le\tfrac32 V.               \tag{BC6}
\]

\paragraph{Normalize the first observation.}
The event $A=\{T\ge1\}$ is decided by the independent seed before any
observation, so its probability $p$ is independent of the latent state.
It is positive: otherwise the expectations at target values zero and
$\rho$ could not both satisfy the bias bound. The finite constant
$c=\mathbb E[Y\mathbf1_{A^c}]$ is also population independent. Under
the seed law conditioned on $A$, form the active procedure
\[
                  Z=pY+c,\qquad T_Z=T.
\]
It preserves every expectation of $Y$ and always takes at least one
observation. The conditioning is implementable by resampling independent
seeds before observing the coin; its constants depend on the specified
policy and known laws, not on $q,D,S$.

For $p<1$, expansion, Cauchy--Schwarz on $A^c$, and completing the square
give
\begin{align*}
 \mathbb E_qZ^2
 &=p\mathbb E_q[Y^2\mathbf1_A]+2c\mathbb E_qY-c^2\\
 &\le pM+2c\mathbb E_qY-c^2/(1-p)\\
 &\le pM+(1-p)(\mathbb E_qY)^2
 \le m:=pM+2\rho^2.
\end{align*}
For $p=1$, take $c=0$ and the final bound holds directly. Apply (BC5) to
$Z$ to obtain $m\ge2\rho/3$ and
$pM\ge2\rho/3-2\rho^2\ge\rho/6$. Set
$P=MW_0$ and $W_{0Z}=\mathbb E_0T_Z=W_0/p\ge1$. Then
\[
 m\le4pM\le4P,\qquad mW_{0Z}\le4P,
 \qquad P\ge\rho/6.                                   \tag{BC7}
\]
Indeed $2\rho^2\le\rho/2\le3pM$. Every work quantity in (BC7) is
all-negative work.

\paragraph{Select good gaps within the capped global class.}
Now use conditional means and moments for $Z$, keeping the same notation.
The bad set $B_-:=\{M_->1/64\}$ has mass at most $64m$. The set
\[
                    B_+:=\{M_+-M_->2m/\rho\}
\]
has mass less than $2\rho$. Otherwise nonatomicity supplies
$E\subset B_+$ of mass $2\rho$, and the admissible
$q=\tfrac12\mathbf1_E$ gives
\[
 \mathbb E_qZ^2
 =\int M_-d\mu+\tfrac12\int_E(M_+-M_-)d\mu>2m,
\]
contradicting its bound. The response bad set
$B_f:=\{f<3/4\}$ also has mass less than $2\rho$. If not, choose such
an $E$ of mass $2\rho$. Its population has target $\rho$, but
\[
 \mathbb E_{\mathbf1_E/2}Z-\mathbb E_0Z
       =\tfrac12\int_E f\,d\mu<3\rho/4,
\]
whereas the bias bounds force at least $31\rho/32$.

Remove these bad sets and the conditional integrability null sets. The
remaining good set $G$ has
$\mu(G^c)\le64m+4\rho\le70m$. On it,
\[
 |\mathbb E_-Z|\le1/8,\qquad \mathbb E_+Z\ge5/8,
 \qquad M_-\le1/64,\qquad M_+\le1/64+2m/\rho.           \tag{BC8}
\]
The use of sets of mass $2\rho$ is necessary for the cap $q\le1/2$;
the argument does not substitute a pointwise bound $q\le\rho$.

\paragraph{Build a common sign test.}
Let $n=\lceil512(1+m/\rho)\rceil\le1282m/\rho$, the inequality using
$m/\rho\ge2/3$. On the same fixed-gap coin, restart the active procedure
$n$ times with independent seeds and consecutive unused flip blocks.
Output plus if their mean exceeds $3/8$. At a fixed latent state the
restarts are iid: condition on each finite stopping prefix, use the
independence of its unused iid suffix, and sum over the countably many
prefixes. Finite conditional expected work ensures almost-sure
termination for both signs at each good gap. No new latent magnitude is
drawn. The same test is used at every good gap. By (BC8), either mean is
at least $1/4$ from the threshold, and Chebyshev bounds either error by
\[
                 \frac{16(1/64+2m/\rho)}n<1/10.
\]
Its negative expected work at gap $d$ is exactly $nw_-(d)$.

\paragraph{Average negative work over separated annuli.}
For $j\ge3$, put
\[
 I_j=(a4^{-j}/2,a4^{-j}],\qquad
 \mu(I_j)=\tfrac34 16^{-j},\qquad
 J=\left\lfloor\log_{16}\frac3{560m}\right\rfloor,
 \qquad K=J-2.
\]
When $J\ge3$, every retained annulus has mass at least $140m$, so
$G$ occupies at least half of it. Choose $d_j\in G\cap I_j$ with
\[
 w_-(d_j)\le\frac4{\mu(I_j)}\int_{I_j}w_-\,d\mu.
\]
This follows by averaging; the integral is positive since the active
procedure always observes at least once. The choices are existential
proof objects, not inputs to the policy. Successive gap ratios exceed
two, every gap is at most $1/256$, and
\[
 \sum_{j=3}^Jd_j^2w_-(d_j)
       \le\frac{16a^2}{3}\int w_-d\mu
       =\frac{16a^2}{3}W_{0Z}.
\]
Lemma~\ref{lem:finite-sign-entropy}, the restart bound, and (BC7) imply,
for $K\ge K_s$,
\[
 c_sK\log K
 \le\sum_{j=3}^Jd_j^2nw_-(d_j)
 \le\frac{16a^2}{3}\frac{1282m}{\rho}W_{0Z}
 \le\frac{30000a^2P}{\rho}.                            \tag{BC9}
\]
The last numerical constant is conservative, since
$(16/3)\cdot1282\cdot4<30000$.

\paragraph{Obtain one positive constant for every rarity.}
For sufficiently small $\rho$, if $P\ge\sqrt\rho$ the desired order
already follows because $\sqrt\rho LH\to0$. Otherwise (BC7) gives
$m\le4\sqrt\rho$ and
\[
 K\ge\frac{\tfrac12\log(1/\rho)-\log(2240/3)}{\log16}-3.
\]
For all sufficiently small universal $\rho$ this exceeds $K_s$, is at
least a positive universal multiple of $L$, and satisfies
$\log K\ge\tfrac12\log L$. Equation (BC9), with the fixed $a=1/4$,
gives $P\ge c\rho LH$. On the remaining compact interval of rarities,
the lower bound $P\ge\rho/6$ in (BC7) allows the same conclusion after
decreasing one positive universal constant. Finally (BC6) transfers the
raw-moment lower bound to $VW_0$, proving (BC1).
\end{proof}

The all-negative cost is the strengthened part of this statement. It
does not interchange separate suprema, assert a same-population product,
or by itself reduce an arbitrary population algorithm to a single coin.
